\section*{Chapitre \ref{ch:mx:fragmentation-and-routing} --- Fragmentation et routage}

\begin{solution}{exo:mx:fragmentation-and-routing:1}
Cours acheteur 10.01 pour 500 (V2~; les 50 actions de V3 forment un lot irrégulier, elles ne sont pas protégées)~; cours vendeur 10.02 pour 1\,600 (V1 200, V3 400, V4 1\,000).
\end{solution}

\begin{solution}{exo:mx:fragmentation-and-routing:2}
\omterm{def:mx:fragmentation-and-routing:locked}{Marché bloqué}, puis \omterm{def:mx:fragmentation-and-routing:locked}{croisé}. Dans le blocage, il n'y a rien à gagner à prix égaux, avant frais~; dans le croisement, un arbitragiste achète sur B à 10.02 et vend sur A à
10.03, soit un centime par action, si les deux cotations sont encore là lorsque ses ordres arrivent.
\end{solution}

\begin{solution}{exo:mx:fragmentation-and-routing:3}
Chaque changement n'est vu qu'avec un retard d'une milliseconde, donc pendant au plus $500\times0.001=50~\%$ du temps. La borne simulée vaut $0.44\times0.0005=0.022~\%$~; la mesure,
0.0098~\%, est plus basse parce que les changements arrivent par rafales dont les fenêtres de retard se chevauchent.
\end{solution}

\begin{solution}{exo:mx:fragmentation-and-routing:4}
$\psi\Omega\psi'=0.76$. Ordre (1, 2)~: $\psi F=(0.8,\,0.346)$, parts 0.842 et 0.158. Ordre (2, 1)~: parts 0.645 pour le marché 2 et 0.355 pour le marché 1.
Bornes~: le marché 1 se situe entre 0.355 et 0.842, le marché 2 entre 0.158 et 0.645.
\end{solution}

\begin{solution}{exo:mx:fragmentation-and-routing:5}
$\alpha_\perp\propto(0.3,0.2)$, puisque $\alpha'\alpha_\perp=-0.06+0.06=0$. \omterm{def:mx:fragmentation-and-routing:cs}{Parts de composante} 0.6 et 0.4~: le marché 1 s'ajuste moins à l'écart et mène la \omterm{def:mx:fragmentation-and-routing:discovery}{formation des prix}.
\end{solution}

\begin{solution}{exo:mx:fragmentation-and-routing:6}
Les \omterm{def:mx:fragmentation-and-routing:cs}{parts de composante} sont des poids normalisés de $\alpha_\perp$, et rien n'oblige les coefficients à partager un même signe. Un poids négatif provient d'un marché dont le
coefficient d'ajustement a le signe d'un marché qui dépasse la cible~: les cotations de C, en retard de dix secondes, bougent après les autres et d'un montant supérieur à l'écart, si bien que la
composante commune se charge contre elles. La \omterm{def:mx:fragmentation-and-routing:is}{part d'information}, qui est une part de variance, ne peut pas être négative.
\end{solution}

\begin{solution}{exo:mx:fragmentation-and-routing:7}
Des prix milieux à la seconde et cinq retards donnent à B entre 0.37 et 0.59, à A entre 0.23 et 0.44, à C entre 0.13 et 0.28~: B arrive en premier. Les retards de 0.2 et 1 seconde sont plus courts que le
temps que mettent les cotations de chaque plateforme pour rejoindre $P^\ast$ (de quelques secondes à quelques dizaines de secondes dans \texttt{firm.tape}), si bien que l'avance se perd dans le bruit propre
aux cotations de chaque plateforme~; avec des retards de 2 et 10 secondes, A est identifiée. Un échantillonnage plus fin n'aide que si les cotations réagissent plus vite que l'intervalle d'échantillonnage.
\end{solution}

\begin{solution}{exo:mx:fragmentation-and-routing:8}
Le volume n'est pas la \omterm{def:mx:fragmentation-and-routing:discovery}{formation des prix}. Dans la simulation, le volume de C est passé de 128\,100 à 230\,900 actions lorsqu'un arbitragiste a commencé à négocier contre ses cotations
périmées, alors que sa \omterm{def:mx:fragmentation-and-routing:is}{part d'information} restait entre 0.01 et 0.04~: le surcroît de volume venait d'autres intervenants qui prenaient les prix tardifs de C. Il faut estimer les \omterm{def:mx:fragmentation-and-routing:is}{parts
d'information} avant et après, et non les parts de volume.
\end{solution}

\begin{solution}{pb:mx:fragmentation-and-routing:1}
\textbf{1.} Un même titre négocié sur plusieurs plateformes à la fois~; une concurrence portant sur les frais, la vitesse et les fonctionnalités, face à une liquidité fragmentée et au routage.
\textbf{2.} Trois flux \texttt{firm.tape} sur un même prix efficient, vus en même temps sur A, avec 2 secondes de retard sur B et 10 sur C, et 70~\% et 45~\% de l'activité de A sur B
et C.
\textbf{3.} 38~\%, 30~\% et 32~\% sur A, B et C, pondérés par le temps.
\textbf{4.} Des coûts de transaction plus faibles et des exécutions plus rapides, davantage de volatilité à court terme, des prix plus proches d'une marche aléatoire.
\textbf{5.} Le meilleur cours de chaque plateforme lui parvient au bout de 0.5 ms~; il calcule le NBBO à partir des cotations en lots standard et le publie 20 microsecondes plus tard.
\textbf{6.} 0.0098~\%~; 0.002~\% à 50 microsecondes, 0.018~\% à 1 ms, 0.084~\% à 5, 0.33~\% à 20, 1.5~\% à 100.
\textbf{7.} 18.5~\% des transactions, presque toutes celles de l'arbitragiste~; 0.02~\% de celles des intervenants de fond.
\textbf{8.} Des désalignements plusieurs fois par seconde sur les titres très actifs, durant une à deux millisecondes.
\textbf{9.} Meilleur cours acheteur d'une plateforme égal (\omterm{def:mx:fragmentation-and-routing:locked}{bloqué}) ou supérieur (\omterm{def:mx:fragmentation-and-routing:locked}{croisé}) au meilleur cours vendeur d'une autre.
\textbf{10.} Blocage 34.4~\% et croisement 10.0~\% de l'heure~; 69 épisodes de croisement, de durée médiane 4.8 secondes.
\textbf{11.} Elle supprime les croisements (565 épisodes, médiane 41 microsecondes, 0.07~\% du temps) et laisse subsister les blocages (42.6~\%, médiane 0.44 seconde).
\textbf{12.} Surtout les cotations périmées de C (volume de C passant de 128\,100 à 230\,900)~: arbitrage de latence.
\textbf{13.} Abroger la Rule 611 et la Rule 610(e), avec une période de commentaires de 60 jours.
\textbf{14.} $\Delta p_t=\alpha\beta'p_{t-1}+\sum\Gamma_k\Delta p_{t-k}+e_t$~; $\mathrm{IS}_j=([\psi F]_j)^2/(\psi\Omega\psi')$.
\textbf{15.} A de 0.63 à 0.79, B de 0.19 à 0.34, C de 0.01 à 0.04~; \omterm{def:mx:fragmentation-and-routing:cs}{parts de composante} 0.75, 0.44 et $-0.20$.
\textbf{16.} Les innovations des plateformes sont corrélées~; la factorisation de Cholesky attribue la part commune au marché qui vient en premier.
\textbf{17.} B en premier (0.37 à 0.59), A en second (0.23 à 0.44)~: les retards sont plus courts que le temps de réaction propre aux cotations et se perdent dans le bruit.
\textbf{18.} \emph{Résultat nommé}~: avec des retards de 0, 2 et 10 secondes, les \omterm{def:mx:fragmentation-and-routing:is}{parts d'information} valent 0.63 à 0.79, 0.19 à 0.34 et 0.01 à 0.04 (\omterm{def:mx:fragmentation-and-routing:cs}{parts de composante}
0.75, 0.44, $-0.20$)~; les meilleurs cours direct et consolidé diffèrent pendant 0.002~\% du temps avec un retard de 50 microsecondes, 0.0098~\% pour une demi-milliseconde,
0.084~\% à 5 ms et 1.5~\% à 100 ms, et dans 18.5~\% des transactions, celles de l'arbitragiste rapide.
\textbf{19.} Non~: le volume de C a presque doublé avec l'arbitragiste alors que sa \omterm{def:mx:fragmentation-and-routing:is}{part d'information} restait proche de zéro.
\textbf{20.} Sur la plateforme dont les intervenants apprennent en premier, et dans les ordres de l'arbitragiste qui le transmettent aux autres.
\end{solution}

\begin{solution}{iq:mx:fragmentation-and-routing:1}
Le meilleur cours acheteur et le meilleur cours vendeur parmi les \omterm{def:mx:fragmentation-and-routing:protected}{cotations protégées} de toutes les plateformes, calculés par le processeur d'information sur les titres à partir des données des plateformes~; un trader
disposant de flux directs les calcule plus tôt et à son propre rythme, si bien que les deux diffèrent chaque fois qu'un changement n'est pas encore parvenu au processeur.
\iqlookfor{\omterm{def:mx:fragmentation-and-routing:protected}{Cotations protégées}~; processeur~; latence.}
\end{solution}

\begin{solution}{iq:mx:fragmentation-and-routing:2}
Un cours acheteur sur une plateforme égal au cours vendeur d'une autre. Les règles américaines interdisent d'afficher des cotations qui bloquent des \omterm{def:mx:fragmentation-and-routing:protected}{cotations protégées} ailleurs~; une plateforme recote ou rejette
l'ordre (ou le route pour prendre la cotation), car un blocage envoie des signaux contradictoires et invite à des boucles de routage.
\iqlookfor{Définition~; la règle~; recotation.}
\end{solution}

\begin{solution}{iq:mx:fragmentation-and-routing:3}
\omterm{def:mx:fragmentation-and-routing:is}{Parts d'information} de Hasbrouck et \omterm{def:mx:fragmentation-and-routing:cs}{parts de composante} de Gonzalo--Granger, issues d'un modèle vectoriel à correction d'erreur appliqué à des cotations synchronisées. Pièges~: des bornes larges
lorsque les innovations sont corrélées (échantillonner plus finement), une avance plus courte que l'intervalle d'échantillonnage, des cotations périmées, et la confusion du volume avec la \omterm{def:mx:fragmentation-and-routing:discovery}{formation des prix}.
\iqlookfor{VECM~; bornes~; fréquence d'échantillonnage.}
\end{solution}

\begin{solution}{iq:mx:fragmentation-and-routing:4}
La plateforme, l'horodatage de la plateforme et mon horodatage de réception, le numéro de séquence, et si la cotation est protégée (lot standard, affichée)~; j'en ai besoin
pour savoir l'ancienneté de chaque prix, détecter les trous et router vers ce qui est réellement disponible.
\iqlookfor{Horodatages~; séquence~; péremption.}
\end{solution}

\begin{solution}{iq:mx:fragmentation-and-routing:5}
Les deux effets existent~: la concurrence fait baisser les frais et les spreads, tandis que la fragmentation de la liquidité augmente les coûts de recherche et de routage. Examiner des titres appariés ou des expériences naturelles
(arrivée d'une nouvelle plateforme, changements de réglementation) et comparer les spreads, la profondeur entre plateformes, la volatilité et l'efficience des prix.
\iqlookfor{Mécanismes des deux côtés~; identification.}
\end{solution}

\begin{solution}{iq:mx:fragmentation-and-routing:6}
Toute personne dont l'ordre ou la cotation est contrôlé ou valorisé à partir de la vue périmée~: des cotations au repos sur des plateformes lentes, capturées par des traders rapides, des ordres exécutables
routés vers des prix qui ont déjà disparu. Mesurer les exécutions et les mark-outs des transactions survenant pendant que les deux vues diffèrent, comme l'ont fait Bartlett et McCrary
avec des horodatages à la microseconde, ainsi que la part des transactions exécutées avec un NBBO consolidé périmé.
\iqlookfor{Qui en pâtit~; une mesure aux instants des transactions.}
\end{solution}
